% Options for packages loaded elsewhere
\PassOptionsToPackage{unicode}{hyperref}
\PassOptionsToPackage{hyphens}{url}
\documentclass[
]{article}
\usepackage{xcolor}
\usepackage[margin=1in]{geometry}
\usepackage{amsmath,amssymb}
\setcounter{secnumdepth}{5}
\usepackage{iftex}
\ifPDFTeX
  \usepackage[T1]{fontenc}
  \usepackage[utf8]{inputenc}
  \usepackage{textcomp} % provide euro and other symbols
\else % if luatex or xetex
  \usepackage{unicode-math} % this also loads fontspec
  \defaultfontfeatures{Scale=MatchLowercase}
  \defaultfontfeatures[\rmfamily]{Ligatures=TeX,Scale=1}
\fi
\usepackage{lmodern}
\ifPDFTeX\else
  % xetex/luatex font selection
\fi
% Use upquote if available, for straight quotes in verbatim environments
\IfFileExists{upquote.sty}{\usepackage{upquote}}{}
\IfFileExists{microtype.sty}{% use microtype if available
  \usepackage[]{microtype}
  \UseMicrotypeSet[protrusion]{basicmath} % disable protrusion for tt fonts
}{}
\makeatletter
\@ifundefined{KOMAClassName}{% if non-KOMA class
  \IfFileExists{parskip.sty}{%
    \usepackage{parskip}
  }{% else
    \setlength{\parindent}{0pt}
    \setlength{\parskip}{6pt plus 2pt minus 1pt}}
}{% if KOMA class
  \KOMAoptions{parskip=half}}
\makeatother
\usepackage{longtable,booktabs,array}
\usepackage{calc} % for calculating minipage widths
% Correct order of tables after \paragraph or \subparagraph
\usepackage{etoolbox}
\makeatletter
\patchcmd\longtable{\par}{\if@noskipsec\mbox{}\fi\par}{}{}
\makeatother
% Allow footnotes in longtable head/foot
\IfFileExists{footnotehyper.sty}{\usepackage{footnotehyper}}{\usepackage{footnote}}
\makesavenoteenv{longtable}
\usepackage{graphicx}
\makeatletter
\newsavebox\pandoc@box
\newcommand*\pandocbounded[1]{% scales image to fit in text height/width
  \sbox\pandoc@box{#1}%
  \Gscale@div\@tempa{\textheight}{\dimexpr\ht\pandoc@box+\dp\pandoc@box\relax}%
  \Gscale@div\@tempb{\linewidth}{\wd\pandoc@box}%
  \ifdim\@tempb\p@<\@tempa\p@\let\@tempa\@tempb\fi% select the smaller of both
  \ifdim\@tempa\p@<\p@\scalebox{\@tempa}{\usebox\pandoc@box}%
  \else\usebox{\pandoc@box}%
  \fi%
}
% Set default figure placement to htbp
\def\fps@figure{htbp}
\makeatother
\setlength{\emergencystretch}{3em} % prevent overfull lines
\providecommand{\tightlist}{%
  \setlength{\itemsep}{0pt}\setlength{\parskip}{0pt}}
\usepackage{bookmark}
\IfFileExists{xurl.sty}{\usepackage{xurl}}{} % add URL line breaks if available
\urlstyle{same}
\hypersetup{
  pdftitle={Tau levels in axonal and soma compartments in BFCN after pro-neurotrophin treatment},
  pdfauthor={Elina Nemtsov},
  hidelinks,
  pdfcreator={LaTeX via pandoc}}

\title{Tau levels in axonal and soma compartments in BFCN after pro-neurotrophin treatment}
\author{Elina Nemtsov}
\date{March 2026}

\begin{document}
\maketitle

{
\setcounter{tocdepth}{2}
\tableofcontents
}
\section*{Abstract}\label{abstract}
\addcontentsline{toc}{section}{Abstract}

Short summary - do last

\section*{Acknowledgements}\label{acknowledgements}
\addcontentsline{toc}{section}{Acknowledgements}

Your acknowledgements here.

\section*{List of Abbreviations}\label{list-of-abbreviations}
\addcontentsline{toc}{section}{List of Abbreviations}

\begin{itemize}
\tightlist
\item
  LC --- Locus Coeruleus
\item
  p75NTR --- p75 Neurotrophin Receptor
\end{itemize}

\section{Background}\label{background}

Placeholder

\subsection{Traumatic Brain Injury and Secondary Retrograde Neurodegeneration}\label{traumatic-brain-injury-and-secondary-retrograde-neurodegeneration}

\subsection{Neurotrophin Signaling and p75NTR in Injury-Induced Degeneration}\label{neurotrophin-signaling-and-p75ntr-in-injury-induced-degeneration}

\subsection{Basal Forebrain Cholinergic Neurons as a Model of Retrograde Vulnerability}\label{basal-forebrain-cholinergic-neurons-as-a-model-of-retrograde-vulnerability}

\subsection{Tau background and roles}\label{tau-background-and-roles}

\subsubsection{Core roles of Tau (physiology)}\label{core-roles-of-tau-physiology}

\subsubsection{When Tau becomes pathological}\label{when-tau-becomes-pathological}

\subsubsection{Post-translational modifications and why phosphorylation matters}\label{post-translational-modifications-and-why-phosphorylation-matters}

\subsubsection{Antibodies used in this study (tau phosphorylation readouts)}\label{antibodies-used-in-this-study-tau-phosphorylation-readouts}

\subsection{proNGF--p75NTR signaling and tau phosphorylation}\label{prongfp75ntr-signaling-and-tau-phosphorylation}

\subsubsection{Linking p75NTR Signaling, Tau, and Axonal Fragmentation After TBI}\label{linking-p75ntr-signaling-tau-and-axonal-fragmentation-after-tbi}

\subsection{Scope and Aims of Thesis}\label{scope-and-aims-of-thesis}

\section{Results}\label{results}

Placeholder

\subsection{Total Tau increase in response to ProNGF treatment}\label{total-tau-increase-in-response-to-prongf-treatment}

\subsubsection{Only somas repeated WB on same samples as above}\label{only-somas-repeated-wb-on-same-samples-as-above}

\section{Materials and Methods}\label{materials-and-methods}

\subsection{Experimental Design}\label{experimental-design}

\subsubsection{Neuronal cultures}\label{neuronal-cultures}

Wild-type pregnant rats were euthanized by exposure to CO2 and soaked in 70\% ethanol for 5 min for sterilization.
Embryonic day 16 (E16) rat fetuses were removed under sterile conditions and kept in phosphate-buffered saline (PBS) on ice. Basal forebrains were dissected and dissociated in serum-free medium (SFM) composed of a 1:1 mixture of Eagle's MEM and Ham's F-12 supplemented with insulin (2.5μg/ml), glucose (6 mg/ml), transferrin (100 μg/ml), putrescine (60 μm), progesterone (20 nm), selenium (30 nm), penicillin (0.5 U/ml), and streptomycin (0.5 μg/ml). The cells were then plated in microfluidic chambers (50) attached to glass coverslips in tissue culture dishes that were precoated overnight with poly-d-lysine (0.2 mg/ml). Two hundred microliters of medium was added to the soma compartment, and 100 μl of medium was maintained in the axon compartments to maintain a
medium volume difference that facilitated the growth of the axons
of the neurons through the microgrooves toward the axon compartment. The cells were maintained with the volume difference
between the soma and axon compartments in SFM supplemented
with 1\% B-27 for 5 days at 37°C to obtain compartmentalized
BFCN cultures that could be treated separately at the axons or
somas. For biochemical experiments, BFCNs were plated on the
top compartment of filter inserts with a pore size of 1 μm, allowing neurites to grow through the filter to the bottom compartment,
providing an axon-enriched compartment for axon- or soma-specific
stimulation (51).

\subsection{Statistical Analysis}\label{statistical-analysis}

\section{Discussion}\label{discussion}

Discussion text.

\section{Conclusion}\label{conclusion}

Conclusion text.

\section*{References}\label{references}
\addcontentsline{toc}{section}{References}

\end{document}
